\sscp{Austronesian incursion/migration(s): \tsc{Seram-Tanimbar-Bomberai}}{15-06-09}
It is not clear where the founders of the \tsc{Seram-Tanimbar-Bomberai} subgroup came from.
There is an areal weakening of *p > \it{f,
h}, {\0} found in many (but not all) languages of STB,
\tsc{Ambon-Seram}, \tsc{Aru}, and \tsc{Timor-Babar}.
But the discontinuities in the rest of their historical phonologies
makes it difficult to link STB with any of them.

In the lexicon, some STB languages share words with languages
strung along the southern Islands,
but not with the rest of central Maluku.
At the same time,
some STB languages share words with other languages in central
Maluku that are not found along the southern string of islands
(see Tables \ref{tab:15-03a} and \ref{tab:15-03b}).

The different nodes of STB all contain conservative features.
Watubela (\tsc{Eastern Islands}) retains PMP *q in all positions.
Yamdena (\tsc{Tanimbar-Bomberai}) retains a number of conservative
features, such as *b = \it{b},
*d = \it{d /{\#}{\gap}},
*p = \it{p}{\#} and retains many trisyllables
and final consonants, among other things.
Kur (\tsc{Teor-Kur}) retains *p = \it{p} in all positions.
Several STB languages also maintain a distinction between *ma-p and *ma-b.

It is not yet clear where the homeland is for STB.
Given the conservative features of Yamdena,
it appears that the \tsc{Tanimbar-Bomberai} node of STB may have
spread from Tanimbar to Bomberai.
Whatever the scenario, the \tsc{Nuclear Tanimbar-Bomberai} node
contains languages in both the Tanimbar Islands (Fordata)
and the Bomberai Peninsula
and attests to a direct link between them.
Heavy on-going contact with Papuan languages in \tsc{East Bomberai}
and \tsc{Bomberai Tip} languages clutters the picture considerably. (See \crf{ch:STB}).

There are also some features shared between STB
and SHWNG that need to be explored further.\footnote{
		One example of the similarity between SHWNG
		and STB is having **natu ‘child,
		small’ (reconstructed by \citealt{BlustTrussel2020} to PEMP).
		Examples from STB are currently only known in languages of the
		Bomberai Peninsula: Arguni, Bedoanas \it{natu}, Erokwanas \it{natuʔ} ‘small’,
		Irarutu \it{ntü}
		‘smallest member of a set/offspring,
		small young’, as well as Kowiai \it{net} ‘small’.}
In \crf{ch:STB}
(\srf{11-01}) we noted that STB is somewhat weakly defined
and it may be better to identify its constituent nodes (\tsc{East
Seram}, \tsc{Eastern Islands}, \tsc{Teor-Kur}, \tsc{Tanimbar-Bomberai})
as independent subgroups.
Given this, and the similarities between STB
and SHWNG, it could be that (parts of?) STB were
settled from “northeast Indonesia” (\srf{15-06-03}).

But, like the problem above with \tsc{Sula-Buru}
and \tsc{Ambon-Seram}, it is difficult to line up the discontinuities
in the sound correspondences to see STB deriving other subgroups
such as \tsc{Sula-Buru}, \tsc{Ambon-Seram}, \tsc{Timor-Babar}, or \tsc{Bima-Lembata}.
Or vice versa.
Again, from the evidence in the languages,
it looks like they got off different boats.